\chapter*{Preface}
Geometry is the oldest part of mathematics that is still full of surprises. Draw any triangle, trisect its
angles, and the trisectors meet in the corners of an equilateral triangle: Euclid could have found this, and
nobody did until Frank Morley in 1899. Put two circles one inside the other and fit a chain of circles between
them, each touching its neighbours: if the chain closes up once, it closes up wherever you start it. Pack
circles into the gaps between circles forever, and if the first four curvatures are whole numbers, all the
infinitely many others are too. Such facts look like coincidences. This book is about the structures that make
them inevitable.

The first four chapters take a fresh look at the plane. Complex numbers turn the classical theorems about
triangles into short computations; groups of transformations explain why there are exactly seven patterns for
a frieze and why no crystal has five-fold symmetry; inversion in a circle turns tangent circles into parallel
lines and hard problems into easy ones; and the projective plane, born from the perspective drawings of
Renaissance painters, gives parallel lines a meeting point and makes points and lines interchangeable. The
next four chapters are algebraic geometry at its most concrete: conics and the theorems of Pascal and
Poncelet, quadratic forms and the straight lines hidden in curved surfaces, plane curves and the theorem of
B\'ezout on how often they meet, and finally the cubic curves, whose points can be added like numbers, the
starting point of the arithmetic of elliptic curves. The last four chapters widen the view again: lattices and
the densest ways to pack balls, from the hexagonal pattern of the plane to the miraculous lattice $E_8$; the
curved worlds of the sphere and of the hyperbolic plane, where the angles of a triangle no longer add up to
$\pi$; convex bodies and their inequalities, with Buffon's needle as a bonus; and curvature, from the four
vertices of every oval to the theorem of Gauss and Bonnet that ties the shape of a surface to its topology.

The book starts from complex numbers, linear algebra and a little group theory, as in Tome~0, Tome~S1 and
the first chapters of Tome~2. Every theorem is proved. A few deeper results, quoted for their beauty, are
labelled as remarks with a precise reference. Each chapter opens with a question to think about, and ``Pause
and think'' boxes ask you to guess before a key result. Exercises are graded \textbf{$\star$} warm-up,
\textbf{$\star\star$} for a good afternoon and \textbf{$\star\star\star$} gems; each has a hint and a complete
solution that starts with \emph{How to find it}, because finding the idea matters more than checking the details.

Geometry is a subject to be seen, so every chapter has pictures, and every picture is drawn from coordinates
computed in Python, so that what you see is exactly true. Each chapter also has a companion notebook that
tests the theorems on random examples, draws the configurations, and lets you search for new ones: random
triangles obeying Morley, Apollonian gaskets with their integer curvatures, Poncelet polygons that close up,
elliptic curves over finite fields, the 240 shortest vectors of $E_8$, and needles dropped on a floor to
measure $\pi$.
